\section{Discussion}
\label{sec:discussion}

\paragraph{Why cheap interventions fail.}
The interventions in Table~\ref{tab:methods} share an implicit assumption: that
a low-resolution model can be made sensitive enough to recover the misses.
Distillation~\cite{kd} aligns a $640$ student to a $1280$ teacher, yet leaves the
tiny-defect classes unchanged---feature alignment cannot reconstruct pixel
evidence that the $640$ input never contained. Patch reinspection does access
more pixels, but cropping removes the global context and scale distribution the
high-resolution detector was trained on, so isolated patches recover few
defects. Both observations point to the same conclusion: for visibility-limited
defects the bottleneck is the input information itself, and only genuine
full-image high resolution supplies it.

\paragraph{Why selective triage helps.}
VBAD does not attempt to circumvent this bottleneck; it decides when to pay for
it. The oracle triage exceeds uniform high resolution because $640$ and $1280$
are complementary: high resolution recovers visibility-limited defects but
slightly harms easy images through scale-distribution shift, so retaining $640$
where it suffices is beneficial. The learned head captures a usable part of this
signal at a fraction of the cost, though a gap to the oracle remains. That gap
is expected: a Type-V defect is by construction one the $640$ detector missed,
so its evidence at $640$ is weak, bounding how well any head conditioned on
$640$ features can anticipate it.

\paragraph{The reliability boundary.}
Perhaps the most consequential finding for deployment is that $42.8\%$ of missed
defects are unrecoverable even at high resolution, concentrated in sub-pixel
classes with scale-sensitive annotations. No resolution policy addresses these;
they define a ceiling on automated recall. An inspection system that reports not
only detections but also \emph{which images lie near this boundary} lets
operators route uncertain cases to manual review---a more honest and practical
behavior than silently optimizing a metric these instances cannot move.

\paragraph{When does VBAD help? A contrasting dataset.}
VBAD is only beneficial when a dataset actually contains visibility-limited
defects. To probe this boundary we apply the same analysis to GC10-DET~\cite{gc10det}, a public
metallic-surface defect benchmark of a different domain (10 classes, 2{,}292
images). Table~\ref{tab:gc10} contrasts the two datasets. On GC10-DET the
low-resolution detector already reaches a high recall ($0.673$), high resolution
yields \emph{no} gain---in fact a slight drop in mAP ($0.657 \to 0.631$)---and
only $8.1\%$ of defects are Type-V, versus $21.5\%$ on the fabric dataset. In
other words, GC10-DET defects are largely visible at low resolution, so there is
little for high-resolution reinspection to recover. The practical consequence is
desirable: a visibility-aware policy triggers high resolution on only a small
fraction of GC10-DET images and therefore neither wastes computation nor incurs
the mild high-resolution degradation. VBAD thus adapts to the dataset---it
contributes where visibility is the bottleneck and stays close to the
low-resolution detector where it is not---rather than applying high resolution
indiscriminately.

\begin{table}[t]
\centering
\caption{When high resolution helps: the fabric dataset is visibility-limited,
the metallic-surface dataset GC10-DET is not. mAP and recall on each test set.}
\label{tab:gc10}
\begin{tabular}{lcc}
\toprule
& Fabric (this work) & GC10-DET \\
\midrule
mAP$_{50}$ @640        & 0.397 & 0.657 \\
mAP$_{50}$ @1280       & 0.482 & 0.631 \\
High-res gain          & $+0.085$ & $-0.026$ \\
Recall @640            & 0.410 & 0.673 \\
Recall @1280           & 0.572 & 0.665 \\
Type-V fraction        & 21.5\% & 8.1\% \\
\bottomrule
\end{tabular}
\end{table}

\paragraph{Limitations.}
Our study uses a single public fabric dataset and two fixed resolutions; the
triage policy is binary rather than a continuous resolution selector, and the
visibility head conditions on global features only. Extending the framework to
multiple resolution tiers, to other surface-inspection domains, and to a
spatially localized visibility map are natural next steps.
